\documentclass{article}
\usepackage[utf8]{inputenc}
\usepackage[T1]{fontenc}
\usepackage{geometry}
\usepackage{graphicx}
\usepackage{float}
\usepackage{hyperref}
\usepackage[english,indonesian]{babel}

\title{Optimasi Pemrosesan Template LaTeX Menggunakan Pendekatan Berbasis AI}
\author{Academic Support Agent}
\date{\today}

\begin{document}
\maketitle

\begin{abstract}

\end{abstract}

\chapter{PENDAHULUAN}


\chapter{TINJAUAN PUSTAKA}


% Original content (commented out):
% % !TeX spellcheck = pt_BR

\documentclass[10pt,a4paper]{article}

\usepackage{fontspec}
\setmainfont{Arial}

\usepackage[margin=2cm]{geometry}

\usepackage{titlesec}

\titleformat{\section}
{\normalfont\fontsize{10}{15}\bfseries}{\thesection.}{2pt}{}
\titleformat{\subsection}
{\normalfont\fontsize{10}{15}\bfseries}{\thesection.}{2pt}{}

\usepackage{graphicx}

\begin{document}
  \begin{center}
    \includegraphics[]{logo.png}\par
    {\bfseries%
    \MakeUppercase{Ministério da Educação -- MEC}\par
    \MakeUppercase{Universidade Federal do Piauí -- UFPI}\par
    \MakeUppercase{Pró-Reitoria de Pesquisa -- PROPESQ}\par
    Coordenadoria Geral de Pesquisa -- CPES\par
    Programa Institucional de Bolsas de Iniciação Científica -- PIBIC\par
    Programa Institucional de Bolsas de Iniciação Científica nas Ações Afirmativas -- PIBIC-AF\par
    Campus Universitário Ministro Petrônio Portela, Bloco 06 -- Bairro Ininga\par}
    {\itshape%
    CEP: 64049-550 -- Teresina-PI -- Brasil -- Fone (86) 3215-5564 -- Fone/Fax (86) 3215-5560\par
    \bfseries%
    E-mail:
    }
    {\bfseries%
    \underline{pesquisa@ufpi.edu.br}\\[1em]
    \MakeUppercase{ANEXO II}\par
    Edital PIBIC e PIBIC nas Ações Afirmativas (AF) / CNPQ e UFPI (2019/2020)\par}
    \underline{Dados do Projeto e do Proponente}
  \end{center}%
  \begin{table}[ht]
    \begin{tabular}{|r|p{12cm}|}
      \hline
      \textbf{\begin{tabular}[c]{@{}r@{}}Proponente\\ (Centro, Departamento)\end{tabular}} & Professor (Centro X, Departamento Y). \\ \hline
      \textbf{\begin{tabular}[c]{@{}r@{}}Título do Projeto\\ \\ Coordenador do\\ Projeto\\ \\ Instituição\end{tabular}} & \begin{tabular}[c]{@{}l@{}}Título.\\ \\ Professor.\\ \\ Projeto vinculado à UFPI, cujo coordenador é o próprio proponente.\end{tabular} \\ \hline
      \textbf{\begin{tabular}[c]{@{}r@{}}Colaboradores\\ (Centro, Departamento)\end{tabular}} & X. \\ \hline
      \textbf{\begin{tabular}[c]{@{}r@{}}Discentes\\ (Curso)\end{tabular}} & Y. \\ \hline
      \textbf{\begin{tabular}[c]{@{}r@{}}N° do Cadastro do\\ Projeto na\\ CPES/PROSPEQ\end{tabular}} & Z. \\ \hline
      \textbf{Grande Área/ Área} & X/ Y. \\ \hline
    \end{tabular}%
  \end{table}
  \textbf{Palavras Chave}: Lorem, Ipsum.
  
  \section{Caracterização do Problema}
  
  \section{Objetivos e Metas}
  
  \section{Metodologia e Estratégia de Ação}
  
  \section{Resultados e Impactos Esperados}
  
  \section{Riscos e Dificuldades}
\end{document}


\end{document}